\Askhsh{14123}{4}
\begin{ylh}{2}
    \item 2.3. Απόλυτη Τιμή Πραγματικού Αριθμού
    \item 3.3. Εξισώσεις 2ου Βαθμού
    \item 4.2. Ανισώσεις 2ου Βαθμού
\end{ylh}
Δίνεται το τριώνυμο $f(x) = x^{2} - \ a x\  - \ (a\  + \ 1)$, $x\  \in \ R$, με παράμετρο $a \in R.$

\begin{alist}
\item Για τις διάφορες τιμές της παραμέτρου $a\ $ να βρείτε το πλήθος των ριζών του τριωνύμου.
\monades{7}
\item Αν είναι $a > - 2$, τότε:
\begin{rlist}
\item Να αποδείξετε ότι οι ρίζες του τριωνύμου είναι οι αριθμοί $- 1$ και $a + 1$.
\monades{4}
\item Να βρείτε την τιμή του α για την οποία το μήκος του διαστήματος λύσεων της ανίσωσης $x^{2} - \ a x\  - \ (a + 1) \leq 0$ είναι ίσο με 2024.
\monades{7}
\item Να βρείτε το πρόσημο του $f\left( \dfrac{a}{2} \right)$.
\monades{7}
\end{rlist}
\end{alist}

